\documentclass[12pt,a4paper,UTF8]{ctexart}
\usepackage[a4paper,top=22mm,bottom=22mm,left=22mm,right=22mm,headheight=15pt]{geometry}
\usepackage{amsmath,amssymb,mathtools,bm}
\usepackage{graphicx}
\usepackage{xcolor}
\usepackage{booktabs,longtable,array}
\usepackage{enumitem}
\usepackage{fancyhdr}
\usepackage{caption}
\usepackage{float}
\usepackage{microtype}
\usepackage{setspace}
\usepackage[hidelinks,unicode]{hyperref}
\usepackage[most]{tcolorbox}

\definecolor{ContestBlue}{HTML}{17365D}
\definecolor{SoftBlue}{HTML}{EEF4FA}
\definecolor{SoftGray}{HTML}{F5F6F7}
\hypersetup{pdftitle={第98套},pdfauthor={物理竞赛教学资料}}

\setstretch{1.24}
\setlength{\parindent}{2em}
\setlength{\parskip}{0.35em}
\setlist{itemsep=0.35em,topsep=0.4em}
\captionsetup{font=small,labelfont=bf,justification=centering}
\setcounter{secnumdepth}{0}

\pagestyle{fancy}
\fancyhf{}
\fancyhead[L]{\small 第98套}
\fancyfoot[C]{\thepage}
\renewcommand{\headrulewidth}{0.4pt}

\ctexset{
  section={format=\Large\bfseries\color{ContestBlue},beforeskip=1.2em,afterskip=0.7em},
  subsection={format=\large\bfseries\color{ContestBlue},beforeskip=1.0em,afterskip=0.5em},
  subsubsection={format=\normalsize\bfseries,beforeskip=0.8em,afterskip=0.35em}
}

\makeatletter
\newcommand{\tightlist}{\setlength{\itemsep}{0pt}\setlength{\parskip}{0pt}}
\newcommand{\maxwidth}{\ifdim\Gin@nat@width>\linewidth\linewidth\else\Gin@nat@width\fi}
\newcommand{\maxheight}{\ifdim\Gin@nat@height>0.72\textheight 0.72\textheight\else\Gin@nat@height\fi}
\newcommand{\pandocbounded}[1]{\begingroup\setkeys{Gin}{width=\maxwidth,height=\maxheight,keepaspectratio}#1\endgroup}
\makeatother

\begin{document}

\begin{titlepage}
  \centering
  \vspace*{58mm}
  {\Huge\bfseries\color{ContestBlue} 第98套\par}
  \vfill
  \vspace*{28mm}
\end{titlepage}

\section{第 1 题}\label{ux7b2c-1-ux9898}

一根质量为 \(m\)、长度为 \(L\) 的均匀细棒，可绕通过中心 \(O\)
的水平轴在竖直面内自由转动。开始时细棒静止于水平位置。质量为 \(m/3\)
的孙鹏（可视为质点）从棒上方 \(h\) 高处由静止下落，在距棒右端 \(L/4\)
处抱住细棒，随即沿棒爬向右端。细棒绕 \(O\) 轴的转动惯量为
\(I_0=mL^2/12\)。

\begin{enumerate}
\def\labelenumi{\arabic{enumi}.}
\tightlist
\item
  若碰撞时间极短，求碰撞后细棒开始转动的角速度大小。
\item
  爬行过程中细棒以恒定角速度转动，求孙鹏相对细棒的爬行速度大小。
\item
  自孙鹏落到棒上时计时，经过 \(t\) 后（\(t\)
  小于细棒转过四分之一周所需时间），求孙鹏到转轴 \(O\) 的距离。
\item
  若要求细棒转到与水平方向成 \(30^\circ\) 时，孙鹏已进入距右端不超过
  \(L/8\) 的区域，求下落高度 \(h\) 的范围。
\item
  在第 4 问的临界情形下，细棒转到 \(30^\circ\) 时孙鹏恰好距右端
  \(L/8\)。此刻质量为 \(2m/3\) 的黄俏以竖直速度 \(v\) 下落，在距棒中心
  \(3L/16\) 处抱住细棒，细棒随即停止转动，求 \(v\)。
\item
  随后孙鹏仍以原速向右端爬。若要求细棒保持静止，黄俏应以多大速度向另一端爬行？
\end{enumerate}

\begin{figure}[H]
\centering
\includegraphics[width=0.843\linewidth,height=\textheight,keepaspectratio,alt={图 1-1　初始碰撞位置、细棒转动方向及 30\^{}\textbackslash circ 位置关系。}]{../figures/question-01.pdf}
\caption*{图 1-1　初始碰撞位置、细棒转动方向及 \(30^\circ\) 位置关系。}
\end{figure}

\newpage

\section{第 2 题}\label{ux7b2c-2-ux9898}

一艘总质量为 \(m_0=1.0\times10^4\,\mathrm{kg}\)
的飞船绕地球沿椭圆轨道运动。某瞬间飞船到地心的距离为地球半径的两倍，即
\(r_0=2R_E\)；飞船速度为
\(v_0=6.5\times10^3\,\mathrm{m/s}\)，速度方向与飞船---地心连线的夹角为
\(60^\circ\)。已知

\[R_E=6.5\times10^6\,\mathrm m,\qquad M=6.0\times10^{24}\,\mathrm{kg},\qquad G=6.67\times10^{-11}\,\mathrm{N\,m^2/kg^2}.\]

\begin{enumerate}
\def\labelenumi{\arabic{enumi}.}
\tightlist
\item
  求飞船---地球系统的机械能以及飞船相对地心的角动量大小。
\item
  求近地点、远地点到地心的距离。
\item
  求飞船的绕行周期，以小时表示。
\item
  现要使飞船从该位置逃离地球引力场。采用最省燃料的方式，使喷气方向与飞船瞬时速度相反。喷气时间可忽略，燃料完全转化为气体，气体相对飞船的喷出速度大小为
  \(v_x=2.0\times10^3\,\mathrm{m/s}\)。求至少消耗的燃料质量。
\end{enumerate}

\begin{figure}[H]
\centering
\includegraphics[width=0.529\linewidth,height=\textheight,keepaspectratio,alt={图 2-1　飞船在椭圆轨道上的瞬时位置、径向与速度方向。}]{../figures/question-02.pdf}
\caption*{图 2-1　飞船在椭圆轨道上的瞬时位置、径向与速度方向。}
\end{figure}

\newpage

\section{第 3 题}\label{ux7b2c-3-ux9898}

\textbf{希尔球与洛希极限}

\begin{enumerate}
\def\labelenumi{\arabic{enumi}.}
\tightlist
\item
  质量为 \(M\) 的行星以角速度 \(\omega\) 自转，求同步卫星轨道半径。
\item
  行星为半径 \(R\)、密度 \(\rho_M\) 的均匀球体；卫星也均匀，密度为
  \(\rho_m\)。当卫星表面自引力与行星产生的潮汐力相等时，卫星趋向碎裂并可能形成行星环。按此模型求洛希极限。
\item
  考虑卫星公转与行星自转同向。为什么轨道高于同步轨道的卫星会逐渐远离行星，而低于同步轨道的卫星会逐渐靠近行星？作定性说明。
\item
  质量为 \(M\) 的行星绕质量为 \(M_0\) 的恒星作半径为 \(r_0\)
  的圆周运动，且
  \(M_0\gg M\)。卫星在行星附近的希尔球内不会被恒星夺走。求希尔球半径的近似表达式。
\item
  水星的希尔球半径 \(R_h\approx1.7\times10^8\,\mathrm m\)，同步轨道半径
  \(R_s\approx2.4\times10^8\,\mathrm m\)，洛希极限
  \(R_r\approx3\times10^6\,\mathrm m\)。判断水星能否长期拥有卫星和星环，并说明原因。
\item
  土星的洛希极限与同步轨道半径很接近。判断土星能否拥有星环，并说明原因。
\end{enumerate}

\begin{figure}[H]
\centering
\includegraphics[width=0.9\linewidth,height=\textheight,keepaspectratio,alt={图 3-1　恒星、行星及行星希尔球的尺度关系。}]{../figures/question-03.pdf}
\caption*{图 3-1　恒星、行星及行星希尔球的尺度关系。}
\end{figure}

\newpage

\section{第 4 题}\label{ux7b2c-4-ux9898}

一根长度为 \(l\)、质量为 \(m\) 的均匀绝缘细棒置于 \(xOy\)
平面内，质心位于原点。两端分别固定点电荷 \(+q\) 和 \(-q\)。空间中存在沿
\(+x\) 方向的匀强电场 \(\bm E_0\)。初态细棒处于稳定平衡，正电荷一端指向
\(+x\) 方向。

\begin{enumerate}
\def\labelenumi{\arabic{enumi}.}
\tightlist
\item
  细棒受到微小转动扰动后，求其小振动周期 \(T_0\)。
\item
  再加入垂直纸面向里的匀强磁场 \(\bm B_0\)。初始给细棒沿 \(+y\)
  方向的平动速度 \(v_0\)，同时使细棒逆时针偏转小角
  \(\theta_0\)，初始角速度为零。将 \(v_0\) 与 \(\theta_0\)
  均视为一阶小量，求细棒此后的平动与转动，结果保留到首阶非零小量。
\end{enumerate}

\begin{figure}[H]
\centering
\includegraphics[width=0.761\linewidth,height=\textheight,keepaspectratio,alt={图 4-1　带电绝缘细棒、电场、磁场与初始扰动的方向约定。}]{../figures/question-04.pdf}
\caption*{图 4-1　带电绝缘细棒、电场、磁场与初始扰动的方向约定。}
\end{figure}

\newpage

\section{第 5 题}\label{ux7b2c-5-ux9898}

空间中存在垂直纸面向里的匀强磁场，磁感应强度为 \(B\)。桌面上有一个半径为
\(3r\)、质量为 \(M\) 的均匀圆盘，盘上同心固定一个半径为 \(r\)
的轻质小圆盘，并在右下方沿半径方向固定两个很小的轻滑轮，使导线绳按图绕过。大圆盘在桌面上纯滚动。空间中另有两个定滑轮，导线绳绕过半径为
\(r\) 的轻质动滑轮，动滑轮下以绝缘绳悬挂质量为 \(m\)
的重物。忽略圆盘从桌面掉落的可能。

\begin{enumerate}
\def\labelenumi{\arabic{enumi}.}
\tightlist
\item
  某时刻重物下落速度为 \(v_0\)，求导线回路中的感应电动势
  \(\mathcal E\)。
\item
  回路总电阻为 \(R\)。当系统加速度为零时，求重物稳定下落的速度 \(v\)。
\end{enumerate}

\begin{figure}[H]
\centering
\includegraphics[width=0.651\linewidth,height=\textheight,keepaspectratio,alt={图 5-1　纯滚动复合圆盘，以及单根导线绕过盘上轻滑轮、两个定滑轮和动滑轮的闭合路径。}]{../figures/question-05.pdf}
\caption*{图
5-1　纯滚动复合圆盘，以及单根导线绕过盘上轻滑轮、两个定滑轮和动滑轮的闭合路径。}
\end{figure}

\newpage

\section{第 6 题}\label{ux7b2c-6-ux9898}

体积为 \(V\) 的均匀电介质，内能为 \(U=CT\)，其中 \(C\)
为常量。电极化率在所研究温区内满足

\[\chi_e=\chi_0+\alpha T,\]

因而极化强度为
\(P=\chi_e E\)。忽略体积变化。外电场作用下，当极化强度改变时，外电场对介质所做元功按题设约定为
\(\delta W_{\rm on}=-EV\,\mathrm dP\)。系统经历以下准静态循环：

\begin{enumerate}
\def\labelenumi{\arabic{enumi}.}
\tightlist
\item
  初态为 \((E_1,T_1)\)；
\item
  在温度 \(T_1\) 下等温减小电场至 \(E_2\)；
\item
  保持 \(E_2\) 不变，使温度降至 \(T_2\)；
\item
  在温度 \(T_2\) 下等温增大电场至 \(E_1\)；
\item
  保持 \(E_1\) 不变，使温度升回 \(T_1\)。
\end{enumerate}

设 \(E_1>E_2>0\)、\(T_1>T_2\)，求此循环的热效率。

\begin{figure}[H]
\centering
\includegraphics[width=0.727\linewidth,height=\textheight,keepaspectratio,alt={图 6-1　电介质循环在 E---T 平面中的四个准静态过程。}]{../figures/question-06.pdf}
\caption*{图 6-1　电介质循环在 \(E\)---\(T\) 平面中的四个准静态过程。}
\end{figure}

\newpage

\section{第 7 题}\label{ux7b2c-7-ux9898}

一束单色激光近乎垂直入射到厚度为 \(a\)、折射率为 \(n\)
的平行玻璃板。空气到玻璃界面的振幅透射、反射系数分别为

\[t=\frac{2}{n+1},\qquad r=\frac{1-n}{1+n};\]

玻璃到空气界面的相应系数为

\[t'=\frac{2n}{n+1},\qquad r'=\frac{n-1}{n+1}.\]

入射光复振幅为 \(u_0\)，真空波长为 \(\lambda_0\)。

\begin{enumerate}
\def\labelenumi{\arabic{enumi}.}
\tightlist
\item
  求直接反射光的复振幅。
\item
  求在玻璃板内共经历三次反射后射出的反射光复振幅。
\item
  求所有反射光相干叠加后的总复振幅。
\item
  求反射光完全相消时玻璃板厚度 \(a\) 的最小正值。
\item
  求光强透射率的最小值 \(T_{\min}\)。
\end{enumerate}

\begin{figure}[H]
\centering
\includegraphics[width=0.873\linewidth,height=\textheight,keepaspectratio,alt={图 7-1　激光在平行玻璃板前、后表面的反射与透射。}]{../figures/question-07.pdf}
\caption*{图 7-1　激光在平行玻璃板前、后表面的反射与透射。}
\end{figure}

\newpage

\section{第 8 题}\label{ux7b2c-8-ux9898}

``观测''通常指同一参考系时刻物体的状态；``看见''则由同一时刻到达相机的光线组成。物体不同部位到相机的距离不同，因此照片对应的各光线一般不是同时发出的。

\begin{enumerate}
\def\labelenumi{\arabic{enumi}.}
\tightlist
\item
  一根静长为 \(l\) 的细杆平行于 \(y\) 轴，并相对地面系 \(S\) 以速度
  \(\beta c\) 沿 \(+x\) 方向运动。相机静止在 \(S\) 系的 \(y\)
  轴上，镜头对准原点且离原点足够远。相机拍到细杆经过原点附近时，照片中两端
  \(A\)、\(B\) 是否重合？画出主视图并标出相应长度。
\item
  将细杆换成静止边长为 \(l\)
  的立方体，其三条共点棱分别平行于坐标轴。画出相机拍到的主视图并标出相应长度。
\item
  将细杆换成静止半径为 \(r\) 的球。证明照片轮廓是半径为 \(r\)
  的圆而不是椭圆，并指出照片最左、最右两点分别对应球面上的哪些点。可用球的共动系中相应半径与
  \(+x\) 方向的夹角表示答案。
\end{enumerate}

\begin{figure}[H]
\centering
\includegraphics[width=0.9\linewidth,height=\textheight,keepaspectratio,alt={图 8-1　运动细杆、立方体和球相对相机的位置关系。}]{../figures/question-08.pdf}
\caption*{图 8-1　运动细杆、立方体和球相对相机的位置关系。}
\end{figure}

\end{document}
